% Generated by Listas/make_statements.py from pratica-lecture-06.md, item 16. Do not edit.
\textbf{Source.} Milgrom \& Roberts (1992), \emph{Economics, Organization and Management}, Ch. 8 Exercises, Mathematical Exercise 5, p.~286 {[}PDF 301{]}. Reproduced verbatim. Answer in the source: none in the book.



\begin{quote}

\begin{enumerate}

\def\labelenumi{\arabic{enumi}.}

\setcounter{enumi}{4}

\tightlist

\item

  Suppose instead that each time the parties interact, there is a probability \(p > 0\) that they will interact again, and, correspondingly, a probability of \((1 - p)\) that this will be their last interaction. Then at any point, the probability that they will interact twice more is \(p^2\), the probability of three more interactions is \(p^3\), and so on. Suppose the parties are interested in the expected value of the total payoffs they receive. Thus, each seeks to maximize the sum of the payoff at the current round, plus \(p\) times the payoff that accrues if they interact again, plus \(p^2\) times the payoff if they should interact a third time, plus \(p^3\) times the payoff from a possible fourth interaction, and so on. In summation notation, the objective of each party \(i\) is to maximize \(\Sigma p^n u_{in}\), where \(n\) runs from 0 (the current round) to infinity and \(u_{in}\) is the payoff to party \(i\) at the corresponding round. Show that so long as the payoff \(V\) to each party when trust is offered and honored is strictly larger than the payoff when trust is not offered (which we set equal to zero in the text), then for any values of the gain \(G\) to abusing trust and loss \(L\) to having trust not be honored there is a value of \(p\) that makes it Nash equilibrium behavior to offer trust and honor it at each round. What threats and punishments can be used in achieving this outcome? {[}Note: for any number \(\Delta\), \(\sum_{n=1}^{\infty} p^n \Delta = p\Delta/(1 - p)\).{]}

\end{enumerate}

\end{quote}

